\section{Symmetric Binary Tree}
\label{sec:trees-symmetric}

\problemstatement{A binary tree is symmetric if it is a mirror image of
itself around its centre: the left subtree is the mirror reflection of the
right subtree. Given the root, determine whether the tree is symmetric.}

\begin{patternbox}
\emph{``Mirror,'' ``symmetric,'' ``reflection''} is the identical-trees
lockstep recursion (Section~\ref{sec:trees-same}) with a twist: instead of
comparing left-with-left and right-with-right, we compare the \emph{outer}
pair with each other and the \emph{inner} pair with each other. Mirror
matching crosses the two children.
\end{patternbox}

\subsection*{Understanding the problem carefully}

Symmetry is a statement about the left subtree versus the right subtree.
Two subtrees are mirror images when their roots are equal \emph{and} the
left one's left child mirrors the right one's right child (the outer edges)
\emph{and} the left one's right child mirrors the right one's left child
(the inner edges). The crossing -- left.left with right.right, left.right
with right.left -- is what encodes ``reflection'' rather than ``equality.''

\begin{ideabox}[Mirror-Compare: Outer with Outer, Inner with Inner]
Define \textsc{Mirror}($a$, $b$): they mirror iff both are null, or both
are non-null with equal values \emph{and} \textsc{Mirror}($a$.left,
$b$.right) \emph{and} \textsc{Mirror}($a$.right, $b$.left). The whole tree
is symmetric iff \textsc{Mirror}(root.left, root.right).
\end{ideabox}

\subsection*{Algorithm}

\begin{enumerate}
  \item If root is null, return true.
  \item \textsc{Mirror}($a$, $b$):
  \begin{enumerate}
    \item If both null, return true; if exactly one null, return false.
    \item If $a.val \ne b.val$, return false.
    \item Return \textsc{Mirror}($a$.left, $b$.right) \textbf{and}
          \textsc{Mirror}($a$.right, $b$.left).
  \end{enumerate}
  \item Return \textsc{Mirror}(root.left, root.right).
\end{enumerate}

\subsection*{Dry run}

\dryrun{Tree: root $1$; left $2$ (children $3,4$); right $2$ (children
$4,3$).}

\begin{itemize}
  \item Mirror($2_L$, $2_R$): values equal. Compare outer
        ($2_L$.left$=3$, $2_R$.right$=3$) and inner ($2_L$.right$=4$,
        $2_R$.left$=4$).
  \item Outer $3,3$: equal, children null $\to$ true. Inner $4,4$: true.
  \item All true $\Rightarrow$ \textbf{symmetric}.
\end{itemize}

\subsection*{C++ implementation}

\begin{lstlisting}
#include <bits/stdc++.h>
using namespace std;

struct TreeNode {
    int val;
    TreeNode* left;
    TreeNode* right;
    TreeNode(int v) : val(v), left(nullptr), right(nullptr) {}
};

bool mirror(TreeNode* a, TreeNode* b) {
    if (a == nullptr && b == nullptr) return true;
    if (a == nullptr || b == nullptr) return false;
    if (a->val != b->val) return false;
    return mirror(a->left, b->right) &&    // outer pair
           mirror(a->right, b->left);      // inner pair
}

bool isSymmetric(TreeNode* root) {
    if (root == nullptr) return true;
    return mirror(root->left, root->right);
}

int main() {
    TreeNode* root = new TreeNode(1);
    root->left = new TreeNode(2);
    root->right = new TreeNode(2);
    root->left->left = new TreeNode(3);
    root->left->right = new TreeNode(4);
    root->right->left = new TreeNode(4);
    root->right->right = new TreeNode(3);

    cout << (isSymmetric(root) ? "symmetric" : "not symmetric") << "\n";
    return 0;
}
\end{lstlisting}

\begin{complexitybox}
\textbf{Time Complexity.} $O(n)$ -- each pair of mirrored nodes is compared
once. \textbf{Space Complexity.} $O(h)$ for the recursion stack.
\end{complexitybox}

\begin{takeawaybox}
Symmetry is equality with the two children crossed. Recognising it as ``the
same-tree recursion, but mirror the child pairing'' means you write it
correctly the first time instead of tangling the four child comparisons.
\end{takeawaybox}
